\documentclass[11pt]{article}
\usepackage[margin=1in]{geometry}
\usepackage{amsmath,amssymb,amsthm,booktabs}
\usepackage[hidelinks]{hyperref}
\title{A nonunimodal Langford coverage sequence\\
\large A counterexample to MathDB problem 320197}
\author{Principia Math in Progress}
\date{September 12, 2026}
\begin{document}
\maketitle

\paragraph{Task 1. Resolve the original conjecture.}
Conjecture 6 of Pan \cite[p.~4]{Pan} is false. The following
counterexample violates both the asserted unimodality and the
asserted maximum value.

A solution to \((2,n)\) is a word of length \(2n\) containing two
copies of every \(k\in\{1,\ldots,n\}\), with exactly \(k\) entries
between the copies of \(k\). If \(\theta_k\) is its first occurrence,
its second occurrence is \(\theta_k+k+1\).
The associated sequence counts inclusive interval coverage:
\[
 \mu_i=\#\{k\in\{1,\ldots,n\}:\theta_k\le i\le\theta_k+k+1\},
 \qquad 1\le i\le2n.
 \tag{1}
\]
This is the interpretation of the source's definition consistent
with its identity \(\sum_i\mu_i=(n^2+5n)/2\).
The conjecture requires an index \(i\le n\) with
\(\mu_i=\lfloor(n+3)/2\rfloor\), and requires the sequence to
increase weakly to that index and then decrease weakly.

Take \(n=7\) and
\[
 S=(7,3,1,6,1,3,4,5,7,2,6,4,2,5).
\]
Every required spacing is verified in the following table.
\[
\begin{array}{c|rrrrrrr}
 k&1&2&3&4&5&6&7\\ \hline
 \theta_k&3&10&2&7&8&4&1\\
 \theta_k+k+1&5&13&6&12&14&11&9\\
 \text{difference}&2&3&4&5&6&7&8
\end{array}
\]
The listed endpoints are exactly the positions \(1,\ldots,14\),
so \(S\) is a solution. It also satisfies the source's convention
for choosing one of a reversed pair: its first \(7\) is at position
\(1\), whereas the reversed word's first \(7\) is at position \(6\).

Counting the intervals in (1) gives
\[
 (\mu_1,\ldots,\mu_{14})
 =(1,2,3,4,4,3,3,4,4,4,4,3,2,1).
 \tag{2}
\]
In particular, \(\mu_5=4>\mu_6=3<\mu_8=4\). A weakly unimodal
sequence cannot decrease and subsequently increase. More explicitly,
if a putative mode is at \(i\le5\), its decreasing tail contradicts
\(\mu_7<\mu_8\); if it is at \(i\ge6\), its increasing prefix
contradicts \(\mu_5>\mu_6\). Thus no mode exists.
Also, \(\max_i\mu_i=4\), while the conjectured value is
\(\lfloor(7+3)/2\rfloor=5\).
Finally, the sum in (2) is \(42=(7^2+5\cdot7)/2\), so the correct
total-weight identity survives. This completes the disproof.

\paragraph{Verification.}
The endpoint table and (2) give a complete finite proof.
The accompanying checker independently computes interval coverage
by direct counting and by endpoint differences, verifies all
spacings and the reversal convention, and enumerates the \(26\)
canonical words for \(n=7\). Four have nonunimodal coverage;
ten have a maximum different from \(5\).

\newpage
\paragraph{Task 2. Characterize unimodal coverage.}
Find structural conditions on the arrangement of first and second
occurrences that characterize the Langford words whose coverage
sequences are unimodal.

\paragraph{Task 3. Determine the possible coverage maxima.}
For every admissible \(n\), determine the minimum and maximum of
\(\max_i\mu_i\) over all Langford words, and characterize arrangements
attaining those bounds.

\paragraph{Task 4. Construct families with many peaks.}
Produce explicit infinite families of Langford words with
nonunimodal coverage, and determine how many alternating rises
and falls can occur as a function of \(n\).

\paragraph{Task 5. Enumerate coverage profiles.}
Develop a recurrence or generating function that counts Langford
words by their coverage sequence, and use it to study the proportion
with unimodal coverage and its asymptotic behavior.

\begin{thebibliography}{9}
\bibitem{Pan}
Z.~Pan, \emph{The Multiple Permutation Problem and Some Conjectures},
\href{https://arxiv.org/abs/1007.4663}{arXiv:1007.4663v1},
July 27, 2010, 16 pp.
Definition and reversal convention: p.~1; \(\theta_k\): p.~3;
Conjecture 6: p.~4.
\end{thebibliography}
\end{document}
